\section{Results}
\label{sec:results}

\subsection{Correctness gates}

\Cref{tab:gates} lists the falsification gates. They are not capability results; they establish
that the implementation matches its mathematical specification and that the state contract
holds. The chunked-forward gate is the important one for long context: the split point is
position 50, inside both a 64-token slot-consolidation block and a 32-token summary block, and
the chunked hidden state still matches the full forward to $1.6\times10^{-6}$.

\begin{table}[t]
\centering
\caption{Falsification gates (\texttt{scripts/gate.py}) and unit tests.}
\label{tab:gates}
\small
\begin{tabularx}{\columnwidth}{@{}>{\raggedright\arraybackslash}X r@{}}
\toprule
Gate & Measured \\
\midrule
Delta rule, chunkwise vs.\ sequential & $\max|\Delta o|=4.17\times10^{-7}$ \\
 & $\max|\Delta S|=3.58\times10^{-7}$ \\
Chunked vs.\ full forward (split at 50) & $\max|\Delta h|=1.55\times10^{-6}$ \\
State bytes at $T=512$ and $T=4096$ & $200{,}704$ (constant) \\
MQAR recall, 400 steps & $0.974$ vs.\ chance $0.062$ \\
Unit tests (\texttt{pytest tests}) & 21 passed \\
\bottomrule
\end{tabularx}
\end{table}

\subsection{Decay spectrum (Gap A)}
\label{sec:results-gap-a}

\Cref{tab:gap-a} and \Cref{fig:gap-a} report delayed recall at four distance buckets. The
comparison is deliberately isolated: the slot memory and summaries are disabled, and the decay
bias is frozen so the comparison measures the initialised prior. When the bias is trainable, a
single-rate model simply learns its own spectrum and the ablation shows nothing; that was the
first, invalid version of this experiment (\Cref{app:ledger}).

\begin{table}[t]
\centering
\caption{Delayed recall by distance (8 bindings, length 320, 300 steps). Frozen decay.}
\label{tab:gap-a}
\small
\begin{tabular}{@{}l c c c c c@{}}
\toprule
Variant & Overall & $<32$ & $32$--$127$ & $128$--$255$ & $\geq256$ \\
\midrule
full & 0.984 & 0.984 & 0.984 & 1.000 & 0.969 \\
spectrum & \textbf{1.000} & 1.000 & 1.000 & 1.000 & \textbf{1.000} \\
single rate & 0.758 & 0.984 & 1.000 & 0.750 & 0.297 \\
\bottomrule
\end{tabular}
\end{table}

\begin{figure}[t]
\centering
\begin{tikzpicture}
\begin{axis}[
  width=\columnwidth, height=5.1cm, ybar,
  ymin=0, ymax=1.08, ylabel={Accuracy}, ylabel style={font=\small},
  symbolic x coords={{$<32$},{$32$--$127$},{$128$--$255$},{$\geq256$}},
  xtick=data, x tick label style={font=\scriptsize}, y tick label style={font=\scriptsize},
  enlarge x limits=0.14,
  legend style={font=\scriptsize, draw=none, at={(0.5,1.03)}, anchor=south, legend columns=3},
  axis lines*=left
]
\addplot[fill=sigmaBlue!65, draw=sigmaBlue] coordinates
  {({$<32$},0.984) ({$32$--$127$},0.984) ({$128$--$255$},1.000) ({$\geq256$},0.969)};
\addplot[fill=sigmaTeal!65, draw=sigmaTeal] coordinates
  {({$<32$},1.000) ({$32$--$127$},1.000) ({$128$--$255$},1.000) ({$\geq256$},1.000)};
\addplot[fill=sigmaRed!55, draw=sigmaRed] coordinates
  {({$<32$},0.984) ({$32$--$127$},1.000) ({$128$--$255$},0.750) ({$\geq256$},0.297)};
\legend{full, spectrum, single rate}
\end{axis}
\end{tikzpicture}
\caption{Recall collapses with distance when all heads share one frozen decay rate. The
log-spaced spectrum holds every bucket; the full model is slightly below the delta-only
spectrum at the longest bucket because the comparison mixes additional operators.}
\label{fig:gap-a}
\end{figure}

The single rate is not broken at short distances: it reaches $0.984$ below 32 tokens and
$1.000$ at $32$--$127$. The failure is specifically long-range, $0.297$ beyond 256 tokens,
against $1.000$ for the spectrum. A single timescale cannot simultaneously forget local
context and preserve a binding across hundreds of tokens; the spectrum can.

\subsection{State expansion and chunk summaries (Gaps B and C)}
\label{sec:results-gap-bc}

\Cref{tab:gap-bc} reports the two comparisons whose benefit is \emph{not} established.

\begin{table}[t]
\centering
\caption{Gap B (dense MQAR, 66 bindings) and Gap C (aggregate count). Both tasks are
seed-noise-limited; see the text.}
\label{tab:gap-bc}
\small
\begin{tabularx}{\columnwidth}{@{}l r r@{}}
\toprule
Variant & Accuracy & Params \\
\midrule
\multicolumn{3}{@{}l}{\emph{Capacity: MQAR, 66 bindings (chance $0.125$)}}\\
\texttt{state\_expand}=1 & \textbf{0.754} & 1{,}694{,}911 \\
\texttt{state\_expand}=2 & 0.555 & 2{,}027{,}231 \\
\midrule
\multicolumn{3}{@{}l}{\emph{Aggregate count (chance $0.111$)}}\\
full (with summaries) & 0.703 & 2{,}080{,}744 \\
no summaries & 0.672 & 1{,}817{,}556 \\
\bottomrule
\end{tabularx}
\end{table}

\begin{figure}[t]
\centering
\begin{tikzpicture}
\begin{groupplot}[
  group style={group size=2 by 1, horizontal sep=1.35cm},
  width=0.49\columnwidth, height=4.6cm, ybar,
  ymin=0, ymax=0.9,
  tick label style={font=\scriptsize},
  ylabel style={font=\small}, x tick label style={font=\scriptsize},
  legend style={font=\scriptsize, draw=none},
  axis lines*=left
]
\nextgroupplot[title={\scriptsize Gap B: 66-binding MQAR}, ylabel={Accuracy},
  symbolic x coords={exp1,exp2}, xtick=data]
\addplot[fill=sigmaTeal!65, draw=sigmaTeal] coordinates {(exp1,0.754) (exp2,0.555)};
\draw[dashed, sigmaRed] (axis cs:exp1,0.125) -- (axis cs:exp2,0.125);
\nextgroupplot[title={\scriptsize Gap C: aggregate count},
  symbolic x coords={full,no-sum}, xtick=data]
\addplot[fill=sigmaOrange!70, draw=sigmaOrange] coordinates {(full,0.703) (no-sum,0.672)};
\draw[dashed, sigmaRed] (axis cs:full,0.111) -- (axis cs:no-sum,0.111);
\end{groupplot}
\end{tikzpicture}
\caption{Neither extra state capacity nor the summary ring shows a reliable benefit at the
tested budget. Dashed lines are chance. These are negative/inconclusive results, not hidden
failures.}
\label{fig:gap-bc}
\end{figure}

\textbf{State expansion.} Doubling the number of independent delta states did not improve
dense MQAR; it reduced accuracy from $0.754$ to $0.555$. The task has 66 bindings, above the
practical rank of one $32\times32$ state per head, so one might expect capacity to help. It
did not at 300 steps. Possible explanations are optimisation (the larger model needs more
steps) or interference between duplicated state matrices; the experiment does not distinguish
them. We therefore do not claim a capacity benefit.

\textbf{Chunk summaries.} The aggregate count task improves by $0.031$ with summaries,
$0.703$ vs.\ $0.672$. An earlier run measured $0.797$ with summaries and $0.891$ without,
\ie the opposite sign. The metric is noise-limited at this budget, and the summary change was
not adopted on the strength of either run. The operator is correct and chunk-invariant; its
quality benefit remains unproven.

\subsection{The gated slot read (Gap D)}
\label{sec:results-gap-d}

\Cref{tab:presence} reports the three-arm presence experiment: the previous cosine read, the
ZCA-whitened read on its own, and the shipped gated mechanism (novelty dead-zone, read
temperature, abstention). All arms use identical data, seeds, optimiser, and 800 steps.

\begin{table}[t]
\centering
\caption{Presence recall with frozen embeddings (16 symbols, 8 present, chance $0.50$).
Means over seeds are shown for readability; individual seeds are the evidence.}
\label{tab:presence}
\small
\begin{tabularx}{\columnwidth}{@{}l c c c c@{}}
\toprule
Arm & Seed 0 & Seed 1 & Seed 2 & Mean \\
\midrule
cosine (previous) & 0.562 & 0.539 & 0.555 & 0.552 \\
ZCA-whitened & 0.648 & 0.531 & 0.539 & 0.573 \\
\textbf{gated} & \textbf{0.977} & \textbf{1.000} & \textbf{0.953} & \textbf{0.977} \\
\bottomrule
\end{tabularx}
\end{table}

\begin{figure*}[t]
\centering
\begin{tikzpicture}
\begin{groupplot}[
  group style={group size=3 by 1, horizontal sep=1.25cm},
  width=0.31\textwidth, height=5.0cm, ybar,
  tick label style={font=\scriptsize}, ylabel style={font=\small},
  axis lines*=left, ymin=0
]
\nextgroupplot[title={\scriptsize presence accuracy}, ylabel={Accuracy}, ymax=1.08,
  symbolic x coords={cos,whiten,gated}, xtick=data,
  xticklabels={cosine,whitened,gated}]
\addplot[fill=sigmaBlue!60, draw=sigmaBlue] coordinates {(cos,0.552) (whiten,0.573) (gated,0.977)};
\draw[dashed, sigmaRed] (axis cs:cos,0.50) -- (axis cs:gated,0.50);
\nextgroupplot[title={\scriptsize read entropy (nats)}, ylabel={Entropy}, ymax=3.0,
  symbolic x coords={cos,whiten,gated}, xtick=data,
  xticklabels={cosine,whitened,gated}]
\addplot[fill=sigmaTeal!60, draw=sigmaTeal] coordinates {(cos,1.2595) (whiten,0.7246) (gated,0.4132)};
\draw[dashed, sigmaRed] (axis cs:cos,2.7726) -- (axis cs:gated,2.7726);
\nextgroupplot[title={\scriptsize memory ablation $\Delta$ (nats)}, ylabel={$\Delta$ CE}, ymax=2.6,
  symbolic x coords={cos,whiten,gated}, xtick=data,
  xticklabels={cosine,whitened,gated}]
\addplot[fill=sigmaOrange!65, draw=sigmaOrange] coordinates {(cos,0.1467) (whiten,0.2173) (gated,2.2551)};
\end{groupplot}
\end{tikzpicture}
\caption{Gap D in one view. Whitening sharpens the read (entropy falls) but does not make it
useful; the gated read both sharpens it and makes it load-bearing, as shown by the ablation
delta. Dashed lines: chance accuracy, uniform entropy $\ln16=2.773$ nats.}
\label{fig:gap-d}
\end{figure*}

\begin{table}[t]
\centering
\caption{Gap D diagnostics, seed 0, 16 slots. Uniform read entropy is $\ln16=2.773$ nats.}
\label{tab:gap-d-diag}
\small
\begin{tabularx}{\columnwidth}{@{}l c c c c c@{}}
\toprule
Arm & Entropy & Top $w$ & Key cos & CE & $\Delta_{\mathrm{abl}}$ \\
\midrule
cosine & 1.260 & 0.608 & 0.023 & 0.663 & $+0.147$ \\
whitened & 0.725 & 0.788 & $-0.044$ & 0.517 & $+0.217$ \\
gated & \textbf{0.413} & \textbf{0.820} & 0.004 & \textbf{0.019} & $\mathbf{+2.255}$ \\
\bottomrule
\end{tabularx}
\end{table}

The gated read is not merely sharper; it is discriminating. Match confidence averages
$0.562$ on present queries and $0.026$ on absent ones, \ie the abstention path separates the
two populations without supervision on the gate. The ablation delta shows the same from the
other side: zeroing the memory now costs $2.26$ nats, against $0.15$ before, so the residual
stream alone can no longer answer the task. ZCA whitening sharpens entropy from $1.26$ to
$0.72$ nats but leaves accuracy at chance, because it amplifies every deviation direction,
including context noise.

\paragraph{No-regression controls.}
The gated mechanism must not break the paths it does not target. \Cref{tab:regression} shows
MQAR (delta-rule path), aggregate recall (summary path), and a 600-step real-data next-token
comparison on the tiny preset. MQAR is unchanged to within seed noise, and is slightly higher
for the gated arm ($+0.002$ and $+0.008$). The aggregate task is seed-noisy but the gated arm
is above the previous behaviour on both seeds. Real-data language modeling moves val CE from
$3.0760$ to $3.0826$ nats, a $0.2\%$ difference inside the run-to-run spread; the training
curves overlap throughout (\Cref{fig:lm-check}).

\begin{table}[t]
\centering
\caption{No-regression controls for the gated read. MQAR: 400 steps, 2 seeds. Aggregate:
600 steps, 2 seeds. LM: tiny preset, 3.06M params, 600 steps, identical stream.}
\label{tab:regression}
\small
\begin{tabularx}{\columnwidth}{@{}l c c@{}}
\toprule
Measure & previous & gated \\
\midrule
MQAR seed 0 & 0.988 & \textbf{0.990} \\
MQAR seed 1 & 0.979 & \textbf{0.987} \\
Aggregate seed 0 & 0.680 & \textbf{0.836} \\
Aggregate seed 1 & 1.000 & 0.859 \\
LM val CE (nats) & \textbf{3.0760} & 3.0826 \\
LM PPL & 21.67 & 21.82 \\
\bottomrule
\end{tabularx}
\end{table}

\begin{figure}[t]
\centering
\begin{tikzpicture}
\begin{axis}[
  width=\columnwidth, height=4.4cm,
  xlabel={optimizer step}, ylabel={training CE (nats)},
  xmin=150, xmax=650, ymin=2.9, ymax=4.45,
  tick label style={font=\scriptsize}, label style={font=\small},
  legend style={font=\scriptsize, draw=none, at={(0.03,0.97)}, anchor=north west},
  axis lines*=left, grid=major, grid style={black!10}
]
\addplot[thick, sigmaBlue, mark=*] coordinates {(200,4.2786) (400,3.4066) (600,3.0300)};
\addplot[thick, sigmaOrange, mark=square*, dashed] coordinates {(200,4.2745) (400,3.4105) (600,3.0431)};
\legend{baseline, gated}
\end{axis}
\end{tikzpicture}
\caption{Real TinyStories next-token training, identical stream. The gated slot read does not
change the loss trajectory. Held-out CE is $3.0760$ (baseline) vs.\ $3.0826$ (gated).}
\label{fig:lm-check}
\end{figure}

\subsection{Language modeling at 21M parameters}
\label{sec:results-lm}

The largest model in this report has 20{,}996{,}090 parameters (6 blocks, width 320, 8 heads,
\texttt{state\_expand}=2, 24 slots, 32 summary slots, no MoE). It was trained on TinyStories
(8K byte-level BPE, 300M-token training split) for 10{,}604 steps in two resumed sessions on
two T4 GPUs, a total of roughly 348M tokens. \Cref{tab:20m} reports the two held-out
evaluations.

\begin{table}[t]
\centering
\caption{21M TinyStories run, 2\,$\times$\,T4, fp16 DDP. Held-out evaluation is 20 batches of
4\,$\,$ $\times$\,1024 tokens (81{,}840 supervised tokens). MQAR is far out of distribution
and is reported for completeness only.}
\label{tab:20m}
\small
\begin{tabular}{@{}l c c c@{}}
\toprule
Session & Steps & Held-out loss & PPL \\
\midrule
1 & 5{,}442 & 1.3225 & 3.753 \\
2 & 10{,}604 & \textbf{1.2523} & \textbf{3.498} \\
\bottomrule
\end{tabular}
\end{table}

The model writes grammatical, on-genre TinyStories prose from a prompt. The checkpoint
predates the gated slot read; its slot memory is the diffuse version audited in
\Cref{tab:slotaudit}, which is exactly the defect Gap D removes. The language model is useful
as a prototype, not as evidence of general language ability: TinyStories is a restricted
synthetic corpus, the tokenizer is 8K byte-level BPE, and no transformer was trained on the
same data for comparison.

\begin{table}[t]
\centering
\caption{Slot-memory audit of the 21M checkpoint (pre-fix). Uniform read entropy is
$\ln24=3.178$ nats. Zeroing the read costs only $+0.035$ nats: the memory acts as a slow
average. This table motivates the gated read.}
\label{tab:slotaudit}
\small
\begin{tabular}{@{}c c c c@{}}
\toprule
Block & Read entropy & Top weight & Cross-slot cos \\
\midrule
0 & 3.133 & 0.065 & 0.465 \\
1 & 3.109 & 0.073 & 0.525 \\
2 & 3.160 & 0.070 & 0.810 \\
3 & 3.159 & 0.055 & 0.715 \\
4 & 3.048 & 0.131 & 0.716 \\
5 & 3.095 & 0.078 & 0.483 \\
\bottomrule
\end{tabular}
\end{table}

\begin{table*}[t]
\centering
\caption{Representative samples from the step-10{,}604 checkpoint (temperature 0.7, top-$k$ 40),
truncated. These show local coherence, not long-horizon recall.}
\label{tab:samples}
\small
\begin{tabularx}{\textwidth}{@{}>{\raggedright\arraybackslash}X@{}}
\toprule
``Once upon a time, there was a little girl named Lily. She loved to play with her toys and
make them into different things. One day, she found a toy that was different from the others.
It was a big, shiny rock that had a pointy edge\ldots'' \\
``The little dog ran into the garden and saw the flowers. He was so happy and he couldn't wait
to show all his friends. But then, something bad happened. A big dog came running and scared
the little dog away\ldots'' \\
``Tom was very sad because he lost his toy. He looked everywhere for it, but it was nowhere to
be found. He started to cry because he missed his toy so much. His mother was worried and
asked him what was wrong\ldots'' \\
\bottomrule
\end{tabularx}
\end{table*}

\subsection{Cost and state}
\label{sec:results-cost}

\Cref{tab:cost} compares SigmaStorm with the size-matched-in-width transformer control from
the experiment ladder. The control matches width, depth, heads, and FFN multiplier, not total
parameter count: SigmaStorm's operators add parameters, so at the same width it has about
$2.3$--$3.0\times$ more parameters and, on CPU, is roughly $2$--$9\times$ slower per forward
pass. There is no compute-efficiency claim here. The architectural advantage that motivates the
design is memory: state is constant in context (\Cref{tab:gates}), while a transformer KV cache
grows linearly. That advantage only matters if long-context training works at scale, which is
not demonstrated in this report.

\begin{table}[t]
\centering
\caption{Forward cost at equal width (CPU, two threads, batch 1). The transformer control
matches width/depth, not parameter count.}
\label{tab:cost}
\small
\begin{tabular}{@{}l r r r r@{}}
\toprule
Config & $\Sigma$ params & $\Sigma$ s & T params & T s \\
\midrule
tiny $T{=}128$ & 2.08M & 0.031 & 0.68M & 0.0036 \\
tiny $T{=}512$ & 2.08M & 0.102 & 0.68M & 0.0147 \\
tiny $T{=}2048$ & 2.08M & 0.411 & 0.68M & 0.1404 \\
20M $T{=}128$ & 20.97M & 0.094 & 8.37M & 0.0180 \\
20M $T{=}512$ & 20.97M & 0.321 & 8.37M & 0.0761 \\
20M $T{=}2048$ & 20.97M & 1.339 & 8.37M & 0.5872 \\
\bottomrule
\end{tabular}
\end{table}

\begin{figure}[t]
\centering
\begin{tikzpicture}
\begin{axis}[
  width=\columnwidth, height=4.4cm,
  xlabel={sequence length $T$}, ylabel={forward time (s)},
  xmode=log, ymode=log, log basis x=2,
  xtick={128,512,2048}, xticklabels={128,512,2048},
  tick label style={font=\scriptsize}, label style={font=\small},
  legend style={font=\scriptsize, draw=none, at={(0.03,0.97)}, anchor=north west},
  axis lines*=left, grid=major, grid style={black!10}
]
\addplot[thick, sigmaBlue, mark=*] coordinates {(128,0.0938) (512,0.3214) (2048,1.3388)};
\addplot[thick, sigmaOrange, mark=square*, dashed] coordinates {(128,0.0180) (512,0.0761) (2048,0.5872)};
\legend{SigmaStorm 21M, transformer control 8.4M}
\end{axis}
\end{tikzpicture}
\caption{Forward cost of the 21M configuration and the width-matched transformer control on
CPU. Both are linear in $T$, but SigmaStorm's operators cost more per token at this
implementation. The relative gap is an engineering problem, not a property of the
architecture.}
\label{fig:cost}
\end{figure}
